% =====================================================================
\section{其他强子—强子过程}\label{sec:dy-other}
% =====================================================================

\cref{thm:dy-factorization}的论证用到了被测末态的三个性质：(i) 被测系统在硬散射中产生，本身是色单态，其衰变产物不参与强相互作用；(ii) 截面对其余末态（旁观者、未观测的喷注、软辐射）完全单举；(iii) 测量不约束旁观者的末态，在割线两侧作为恒等算符出现（\cref{rem:dy-one-time-per-jet}）。凡具有这三个性质的过程都继承 DY 的地位；被测末态带色的过程在 (iii) 处失效，其因子化至今只是猜想。下面分别陈述。

\begin{theorem}[色单态末态 \tagA]\label{thm:dy-colour-singlet}
设被测的只是一个在硬散射中产生的色单态系统 $V$：$W^\pm$、$Z$（含其轻子衰变）；Higgs 玻色子（经胶子融合，用重顶夸克极限下的有效顶点或保留顶夸克圈的完整理论，或经其他色单态道产生），其衰变产物不带色；电弱规范玻色子对 $W^+W^-$、$ZZ$、$W^\pm Z$（轻子衰变）。对其余末态单举，则 $V$ 的质量与快度分布具有\cref{thm:dy-factorization}的形式：部分子分布与 DY、DIS 中的相同，只是硬函数 $\mathcal H_{ij}$ 换成相应过程的硬函数，余项同样是 $\order{(\Lambda/Q)^p}$，$Q$ 为硬标度。\\
非微扰地位：\cref{asm:intro-bridge}（MA1、MA2）。
% errtype: E-S
\uses{thm:dy-factorization,prop:dy-regions,lem:dy-fsi,prop:dy-deformation,def:pdf-gluon,prop:appI-larin,asm:intro-bridge}
\end{theorem}

\begin{proofsketch}
(1) \cref{thm:dy-factorization}的证明中，虚光子只以“在硬顶点 $V_H$ 处产生的色单态系统”出现：区域分析、割图求和、单侧改写与形变（\cref{prop:dy-regions,lem:dy-fsi,prop:dy-deformation}）都只涉及两个强子的共线子图、软子图与 Glauber 交换，不涉及 $V$ 本身；所测的 $V$ 的变量不约束强子末态。(2) 对 $W$、$Z$ 与规范玻色子对，改变的只是 $H$ 中的电弱顶点；对无质量夸克，$\order{\as}$ 的 QCD 修正与虚光子相同，只需替换电弱耦合（维数正规化中轴矢流的 $\gamma_5$ 按\cref{prop:appI-larin}（\tagR，$R_{\mathrm{PT}}$）处理）。(3) 对胶子融合产生 Higgs，活动部分子是胶子，共线因子是胶子分布（\cref{def:pdf-gluon}），其 Wilson 线在伴随表示中；\cref{app:appO}的分析只用到活动部分子无质量、强子为色单态，与活动部分子是夸克还是胶子无关。重顶夸克有效顶点的匹配是一个额外的、逐阶可算的硬步骤。完整论证与 DY 相同\cite[第~14.4~节]{Collins2011}；硬函数见\cite{AltarelliEllisMartinelli1979,HambergVanNeervenMatsuura1991}（DY 与 $W$、$Z$）及\cite{Dawson1991,DjouadiSpiraZerwas1991,HarlanderKilgore2002,AnastasiouEtAl2015}（胶子融合，NLO 到 N$^3$LO）。
\end{proofsketch}

\paragraph{地位与缺口}\tagA。与\cref{thm:dy-factorization}完全相同：Glauber 区的输入\cref{asm:appO-mixing}(c)、\cref{asm:appO-deformation}(b)--(d)、\cref{asm:appO-rewritten-cs}等，附录~N 的一般输入，具名缺口\cref{asm:appN-brst-nonlocal}（\tagC，见\cref{thm:method-template}），余项幂次 $p$ 不由证明确定（\tagC）。没有额外的非微扰输入，但要求 $V$ 的衰变产物不带色；若测量带色的衰变产物（例如 $H\to b\bar b$ 的 $b$ 喷注），就落入\cref{conj:dy-coloured-final}的情形。对 $V$ 的横动量 $q_T\ll Q$ 的微分截面，适用的是 TMD 因子化（\cref{thm:tmd-factorization}），软因子不再为 $1$。

\begin{theorem}[极化 Drell--Yan \tagA]\label{thm:dy-polarized}
对纵向极化（螺旋度）与横向极化（横自旋）强子的 DY，自旋相关的截面差（例如两个强子螺旋度相同与相反之差）具有\cref{thm:dy-factorization}的因子化形式，其中的部分子分布换成螺旋度分布与横度分布（\cref{def:pdf-polarized}），硬函数换成相应的自旋投影硬函数。\\
非微扰地位：\cref{asm:intro-bridge}（MA1、MA2、MA4）。
% errtype: E-S
\uses{thm:dy-factorization,def:pdf-polarized,def:appI-hvbm,asm:intro-bridge}
\end{theorem}

\begin{proofsketch}
把\cref{thm:dy-factorization}的证明逐步重复，只在共线因子与硬部分的连接处插入自旋投影：纵向极化取 $\gamma^+\gamma_5$ 投影，得到螺旋度分布；横向极化取 $\gamma^+\gamma^j_\perp\gamma_5$ 型投影，得到横度分布。Glauber 区与软胶子的处理与自旋无关，因为领头幂次的软与 Glauber 耦合只经过共线线的大分量，不改变其自旋结构。领头阶时，矢量流保持手征，入射夸克与反夸克的螺旋度相反，部分子层次的双螺旋度不对称为 $-1$。横度分布是手征奇的，在 DIS 中不出现于领头扭度，在两个横向极化强子的 DY 中则以两个手征奇分布之积的形式出现\cite{RalstonSoper1979,JaffeJi1991}。极化硬函数的计算需要维数正规化中 $\gamma_5$ 的处理（\cref{def:appI-hvbm}）。
\end{proofsketch}

\paragraph{地位与缺口}\tagA。与\cref{thm:dy-factorization}相同；此外极化部分子分布的 $\MSbar$ 定义依赖 $\gamma_5$ 方案（HVBM 加 Larin 有限重整化），这是方案约定而不是缺口。

\begin{conjecture}[强子对撞中有色末态的共线因子化 \tagC]\label{conj:dy-coloured-final}
单举的大横动量强子与喷注以及重夸克对的强子—强子产生满足标准的共线因子化。例如，对 $A+B\to h(p_h)+X$（$p_{h,T}\sim\sqrt s$），
\[
E_h\frac{\dd\sigma}{\dd^3p_h}=\sum_{i,j,k}\int\dd\xi_A\,\dd\xi_B\,\frac{\dd z}{z^2}\,f_{i/A}(\xi_A;\mu)\,f_{j/B}(\xi_B;\mu)\,d_{h/k}(z;\mu)\,E_k\frac{\dd\hat\sigma_{ij\to k}}{\dd^3p_k}+\order{\Big(\frac{\Lambda}{p_{h,T}}\Big)^{p}},
\]
其中 $d_{h/k}$ 是碎裂函数（\cref{def:pdf-ff}），$p_h=zp_k$，因子 $z^{-2}$ 来自 $\dd^3p_h/E_h=z^2\,\dd^3p_k/E_k$；对红外安全的喷注截面与重夸克对总截面，类似地有 $\sigma=\sum_{ij}f_i\otimes f_j\otimes\hat\sigma_{ij}$。
\uses{thm:dy-factorization,def:pdf-ff}
\end{conjecture}

\paragraph{地位与缺口}\tagC。这些公式被普遍使用，但没有与 DY 同等细节的全阶证明。Collins 的教科书写道，这些情形中存在“relative lack of detailed proofs”，并把 Nayak、Qiu 与 Sterman 的分析\cite{NayakQiuSterman2005b}列为例外\cite[第~14.7~节，第~570--571~页]{Collins2011}；Diehl 等人指出，对末态带色的单次散射过程，“to the best of our knowledge the absence of Glauber gluon effects has not yet been proven”\cite[第~5~节]{Diehl:2015bca}。对重夸克产生，CSS 从低阶图论证了标准公式，并强调需要所有阶的分析才能可靠地确立这一猜想\cite{Collins:1985gm}（据其摘要）；Aybat 与 Sterman 给出了有色末态软与 Glauber 相消的纲要\cite{AybatSterman2009}；Catani、de~Florian 与 Rodrigo 发现类空共线极限中严格共线因子化在振幅层次被破坏，并指出它可能从 N$^3$LO 起影响质量奇异性的因子化\cite[第~8~节]{CataniDeFlorianRodrigo2012}。在\cref{app:appO}的框架中，论证在单侧改写一步（(O-3)）失效：被测的有色末态使测量投影不与光锥演化对易，相关的钳夹区是相对于方向对 $(n_A,w_1)$ 的 Glauber 区（\cref{prop:break-measurement}）。另一方面，NLO 的重夸克\cite{NasonDawsonEllis1988}与喷注\cite{Kang:2016mcy}等固定阶计算都以这种形式组织，与数据的符合是证据而不是证明；对同一类过程的 TMD 版本（背对背强子），因子化不成立（第~18~章\cref{thm:break-tmd}），但共线的、对 $q_T$ 积分的截面没有已知的反例；类空共线极限中严格共线因子化的失效（\cref{prop:break-spacelike}）与超出 NNLO 的质量奇异性因子化（\cref{prop:break-mass-singularity}）也在第~18~章讨论。要把本猜想提升为 \tagA，需要对有色末态完成 (O-3) 的替代论证。

\begin{conjecture}[极化强子对撞中的 $A_{LL}$（喷注、$\pi^0$） \tagC]\label{conj:dy-polarized-jets}
纵向极化强子对撞中喷注与 $\pi^0$ 产生的双自旋不对称 $A_{LL}$ 以螺旋度分布按\cref{conj:dy-coloured-final}的形式因子化。
\uses{conj:dy-coloured-final,thm:dy-polarized}
\end{conjecture}

\paragraph{地位与缺口}\tagC。继承\cref{conj:dy-coloured-final}：末态带色，Glauber 相消没有证明；自旋投影本身不引入新的困难（\cref{thm:dy-polarized}）。
